\section{经典模态综合法的数学基础}

给定一个离散化为 $n$ 个自由度的区域 $\Omega$，其刚度矩阵为 $\mathbf{K} \in \mathbb{R}^{n\times n}$，质量矩阵为 $\mathbf{M} \in \mathbb{R}^{n \times n}$。广义特征值问题
\begin{equation} \mathbf{K} \mathbf{u}_l= \lambda_l \mathbf{M} \mathbf{u}_l,
  \label{eq:generalize-eigen-problem}
\end{equation}
给出特征值 $\lambda_l \in \mathbb{R}$ 和特征向量 $\mathbf{u}_l\in\mathbb{R}^{n}$ 的配对。模态分析通常只求解低频模态（即对应于较小 $\lambda_l$ 的模态）。

\subsection{划分与子结构}

要应用模态综合法，需要将区域划分为子结构。区域 $\Omega$ 的一个划分 $\mathcal{P}$ 是一组互不重叠的子结构 $\Omega_i, i=1, \cdots m$，满足 $\bigcup_i\Omega_i = \Omega$（示例见~\cref{fig:domain-partition}）。子结构之间的界面定义为 $\Gamma_{ij} = \Omega_i \cap \Omega_j$。划分 $\mathcal{P}$ 的界面定义为其并集 $\partial \mathcal{P} = \bigcup_{i,j} \Gamma_{ij}$。第 $i$ 个子结构的界面定义为 $\partial \Omega_i = \bigcup_{j} \Gamma_{ij}$。注意符号 $\partial$ 不是指外部区域边界。每个子结构 $\Omega_i$ 内部的自由度记为 $u_i \in \Omega_i\setminus\partial \Omega_i$（不包含界面），所有剩余自由度收集到向量 $u_b$ 中。不失一般性，我们假设自由度已重新排序，使得 $\mathbf{u} = [u_1^T, u_2^T, \cdots, u_m^T, u_b^T]^T$。此时 $\mathbf{K}$ 具有以下分块结构：
\begin{equation}
  \left[
    \begin{array}{ccc|c}
      \mathbf{K}_{11} & & & \mathbf{K}_{1b}\\
      &\ddots & & \vdots\\
      & & \mathbf{K}_{mm} & \mathbf{K}_{mb}\\\hline
      \mathbf{K}_{b1} &\cdots
    &\mathbf{K}_{bm} &\mathbf{K}_{bb}
    \end{array}
  \right]
  \triangleq
  \begin{bmatrix}
    \mathbf{K}_{\mathcal{I}\mathcal{I}}&\mathbf{K}_{\mathcal{I}b}\\
    \mathbf{K}_{b\mathcal{I}}&\mathbf{K}_{bb}
  \end{bmatrix}.
\end{equation}
质量矩阵 $\mathbf{M}$ 具有相同的结构。

\begin{figure}[tbp]
  \centering
  \newcommand{\PIC}[1]{\includegraphics[height=0.28\linewidth]{domain_definition/#1}}
  \PIC{domain-def.png}
  \PIC{inner_modes.png}
  \PIC{interface_modes.png}
  \caption{\textbf{划分、子结构特征模态和界面模态示例。} 区域 $\Omega$ 被划分为四个子结构（左图）。绘制了第 1 个子结构特征模态（中图）和一个界面模态（右图）。}
  \label{fig:domain-partition}
\end{figure}

\subsection{Craig-Bampton (CB) 方法}

本节简要介绍模态综合法的代表性实现——Craig-Bampton 方法（CB 方法），以此概述典型模态综合法的算法流程。

\paragraph{子结构特征模态}
该方法首先并行计算每个子结构的 $q_\mathcal{I}$ 个（截断后的）子结构特征模态 $\mathbf{S}_\mathcal{I}$。其求解方程如下：
\begin{equation}
  \mathbf{K}_{ii} \mathbf{S}_i = \mathbf{M}_{ii} \mathbf{S}_i \Lambda_i.
  \label{eq:sub-eig}
\end{equation}
$\mathbf{K}_{ii}$ 和 $\mathbf{M}_{ii}$ 是全局矩阵 $\mathbf{K}$ 和 $\mathbf{M}$ 的子矩阵，与由 $\Omega_i$ 中单元组装而成的刚度和质量矩阵不同。上述方程实质上描述了一个约束模态分析问题：固定界面 $\partial \Omega_i$ 上的自由度为零。

\paragraph{界面模态}
\label{sec:interface}
随后计算界面模态。这些模态由两部分组成：界面处的单位位移激励和子结构内部的响应，即 $\mathbf{S}_b=[\mathbf{S}_{b\mathcal{I}}^T, \mathbf{S}_{bb}^T]^T$。$\mathbf{S}_{bb} = \mathbf{I}$ 表示界面处的 Dirichlet 边界条件，而
\begin{equation}
  \mathbf{S}_{b\mathcal{I}} = - \mathbf{K}_{\mathcal{I} \mathcal{I}}^{-1}
  \mathbf{K}_{\mathcal{I}b} \mathbf{S}_{bb}= - \mathbf{K}_{\mathcal{I} \mathcal{I}}^{-1}
  \mathbf{K}_{\mathcal{I}b} \mathbf{I},
  \label{eq:traditional-response}
\end{equation}
是该边界条件下的静力平衡状态。由于 $\mathbf{K}_{\mathcal{I}\mathcal{I}}$ 是分块对角矩阵，上述方程同样支持在各个子结构上并行求解。

\paragraph{缩减}
\label{sec:reducing}
计算得到所有子结构特征模态 $\mathbf{S}_i$ 和界面模态 $\mathbf{S}_b$ 后（示意图见~\cref{fig:domain-partition}），可将它们组装为缩减矩阵 $\mathbf{S}$。随后，原始自由度 $\mathbf{u}$ 通过 $\mathbf{S}$ 缩减为 $\mathbf{z}$：
\begin{equation}
  \mathbf{u} =
  \left[
    \begin{array}{ccc|c}
  \mathbf{S}_1 &&&\\
      &\ddots && \mathbf{S}_{b\mathcal{I}}\\
      &&\mathbf{S}_m &\\\hline
      &&&\mathbf{S}_{bb}
    \end{array}
  \right]
  \begin{bmatrix}
  z_1 \\\vdots\\ z_m\\z_b
  \end{bmatrix}
  \triangleq
  \begin{bmatrix}
    \mathbf{S}_{\mathcal{I}} & \mathbf{S}_{b}
  \end{bmatrix}\mathbf{z}
  \triangleq \mathbf{S} \mathbf{z}.
  \label{eq:redmat}
\end{equation}

应用 Galerkin 三重矩阵乘积后，原问题缩减为一个自由度更少的特征值问题：
\begin{equation}
  \mathbf{S}^T\mathbf{K} \mathbf{S} \mathbf{z}_l =\lambda_l \mathbf{S}^T \mathbf{M} \mathbf{S} \mathbf{z}_l.
  \label{eq:reduced-matrix-pencil}
\end{equation}
该问题可以高效求解。

在计算具有以下结构的缩减刚度矩阵和质量矩阵时：
\begin{equation}
  \begin{bmatrix}
    \Lambda_{1} &&&&\\
    &\ddots&&&\\
    &&\Lambda_{m}&&\\
    &&&&\tilde{\mathbf{K}}_{bb}
  \end{bmatrix},
  \begin{bmatrix}
    \mathbf{I}_{1} & & & \tilde{\mathbf{M}}_{1b}\\
    &\ddots & & \vdots\\
    & & \mathbf{I}_{m} & \tilde{\mathbf{M}}_{mb}\\
    \tilde{\mathbf{M}}_{b1} &\cdots &\tilde{\mathbf{M}}_{bm} &\tilde{\mathbf{M}}_{bb}
  \end{bmatrix},
\end{equation}
会涉及大量的大型矩阵-矩阵乘法：
\begin{equation}
  \begin{split}
    \tilde{\mathbf{K}}_{bb} &= \mathbf{K}_{bb} - \sum_i{\textcolor{blue}{\mathbf{K}_{bi}}
      \textcolor{red}{\mathbf{K}_{ii}^{-1} \mathbf{K}_{ib}}}\\
    \tilde{\mathbf{M}}_{ib} &= \textcolor{blue}{\mathbf{M}_{ib} -
    \mathbf{M}_{ii}} \textcolor{red}{\mathbf{K}_{ii}^{-1}\mathbf{K}_{ib}}\\
    \tilde{\mathbf{M}}_{bb} &= \mathbf{M}_{bb} - \sum_i{\color{blue} (\textcolor{red}{\mathbf{K}_{bi}
\mathbf{K}_{ii}^{-1}} \mathbf{M}_{ib} + \mathbf{M}_{bi} \textcolor{red}{\mathbf{K}_{ii}^{-1}
\mathbf{K}_{ib}})}\\
&\quad\quad\quad+\sum_i \textcolor{red}
    {\mathbf{K}_{bi}\mathbf{K}_{ii}^{-1}}\textcolor{blue}{\mathbf{M}_{ii}}
\textcolor{red}{\mathbf{K}_{ii}^{-1} \mathbf{K}_{ib}}.
  \end{split}
  \label{eq:reducing}
\end{equation}
尽管像 $\mathbf{K}_{ii}^{-1}\mathbf{K}_{ib}$ 这样的项已经在 $\mathbf{S}_{b\mathcal{I}}$ 中被预计算，且许多矩阵乘法（彩色文本）可以并行执行，但其计算代价依然高昂。事实上，在某些文献~\cite{aristidou2013dynamic}中，界面模态计算（\Cref{eq:traditional-response}）和缩减过程（\Cref{eq:reducing}）统称为舒尔补。

\subsection{计算代价分析}

\begin{figure}[H]
  \centering
  \includegraphics[width=\linewidth]{substructure_analysis/schur_vs_eig_timing_curves.pdf}
  \caption{\textbf{传统模态综合法并行部分的计算代价分析。} 在一个正方形网格上，以 $N\times N$ 个顶点离散。左图中，$t_{\text{sub-eig}}$ 随特征模态数量呈超线性增长；当子结构特征模态数量较少时，固定代价 $t_{\text{sub-schur}}$（红色虚线）与 $t_{\text{sub-eig}}$ 接近。}
  \label{fig:breakdown-CMS}
\end{figure}

总体代价主要由两部分组成：
\begin{itemize}
    \item $t_{\text{sub-eig}}$：求解子结构特征值问题（\cref{eq:sub-eig}）。
    \item $t_{\text{sub-schur}}$：\cref{eq:traditional-response}
      需要对 $\mathbf{K}_{ii}$ 进行一次分解，随后进行 $n_{b,i}$ 次前代/回代，其中 $n_{b,i}$ 是 $\partial \Omega_i$ 的自由度数。接着是缩减过程中子结构级别的矩阵--矩阵乘法。
\end{itemize}

\Cref{fig:breakdown-CMS}（左）比较了 $t_{\text{sub-eig}}$（实线）与 $t_{\text{sub-schur}}$（虚线）。从图中可以看到，$t_{\text{sub-eig}}$ 会随着所求子结构特征模态数量的增加而呈现明显的超线性增长。这意味着，当需要通过继续增加高频内部模态来改善逼近能力时，特征值求解本身就会迅速成为新的计算瓶颈。相反，$t_{\text{sub-schur}}$ 在模态数量较少时已经表现为一项不可忽略的固定开销；也就是说，即便目标模态数不多，传统模态综合法也必须首先承担一笔较高的界面处理成本。

进一步看，$t_{\text{sub-schur}}$ 虽然表面上包含矩阵分解、前代/回代和矩阵乘法三个步骤，但真正占主导的并不是对 $\mathbf{K}_{ii}$ 的初始分解，而是随后大量重复执行的前代/回代以及缩减阶段的矩阵--矩阵乘法。前者需要针对界面上的每个自由度重复求解线性系统，后者则需要在子结构层面对多个中间量进行装配与乘积计算。因此，舒尔补的代价并不是一次性的，而是会随着界面规模持续放大。

这一点对于复杂几何尤其重要。由于前代/回代的次数与界面自由度数 $n_{b,i}$ 直接相关，而矩阵乘法的规模也受界面维度影响，随着子结构之间的耦合边界变得更加复杂，$t_{\text{sub-schur}}$ 会迅速增长。换言之，界面自由度不仅决定了界面模态构造的规模，也直接决定了界面缩减步骤的时间开销。对于界面较大、拓扑较复杂或划分边界较曲折的模型，这一部分代价往往比直观上预期得更为显著。

除了时间代价之外，舒尔补还伴随着明显的内存负担。其缩减过程中涉及的界面矩阵通常趋于稠密，相关中间量的生成、存储与后续运算都会占用大量内存。因此，在超大规模问题中，舒尔补不仅会拖慢计算速度，还可能先于其他步骤触及内存上限。这也是为什么传统模态综合法在大规模并行环境下常常并非首先受限于浮点运算能力，而是受限于界面处理过程中的存储与数据搬运成本。

综合来看，传统模态综合法实际上同时受到两类瓶颈的制约：一方面，界面处理带来的舒尔补具有较高且难以回避的固定开销，并且会随着界面自由度数目增长而进一步恶化；另一方面，为弥补表达不足而继续增加高频内部模态，又会导致 $t_{\text{sub-eig}}$ 的超线性增长。前者使方法在界面处理上“起步即昂贵”，后者则使方法在提高精度时“越算越慢”。这两类代价叠加之后，便构成了传统模态综合法在超大规模特征值问题中的核心计算障碍。

基于上述分析，后文方法的设计目标也因此变得明确：一方面，希望用尽可能少但更有效的子结构特征模态来降低 $t_{\text{sub-eig}}$ 的增长压力；另一方面，希望绕开计算与内存均十分昂贵的舒尔补步骤，从而降低界面处理的总体代价。换言之，后续方法并不是单纯在现有流程上做局部优化，而是试图同时针对这两类瓶颈重新组织降维特征子空间与界面模态的构造方式。
